\section{Introduction}
\label{sec:intro}

The rapid proliferation of containerized workloads has transformed the way computational resources are packaged, deployed, and shared across modern infrastructure.
In particular we are talking about Containers, but what are they?\\
\textbf{Containers} are lightweight and portable execution environments built upon Linux kernel primitives such as namespaces, control groups, and security computer facilities, and in the last years have become the main deployment unit for cloud-native applications, data processing pipelines, artificial intelligence and machine learning workloads \cite{sultan2019container}.
The emergence of large-scale AI training and inference has introduced a new and critical dimension to this landscape: the need to share hardware accelerators, most prominently \textbf{NVIDIA Graphics Processing Units (GPUs)}, across multiple isolated container workloads.
Unlike CPU resources, which are abstracted transparently by the kernel scheduler, GPU access requires deep integration between the container runtime and the host operating system's driver stack \cite{nvidia_ctk_docs}. This integration is precisely where the attack surface examined in this paper resides.
Going into details: The \textbf{NVIDIA Container Toolkit} (\texttt{nvidia-ctk}) is the software layer responsible for exposing GPU capabilities to containerised processes \cite{nvidia_ctk_github}.
It acts as a privileged extension of the container runtime, executing hooks with root privileges on the host to configure the necessary device mappings and library bindings before the container process starts; the negative side of things is that it can be fatally exploited when the privileged process fails to sanitise the environment it inherits from an untrusted source.
This vulnerability is identified as \textbf{CVE-2025-23266} \cite{nvd_cve_2025_23266} and allows a process running inside a container (potentially an unprivileged tenant in a shared cloud environment) to inject the \texttt{LD\_PRELOAD} environment variable into the execution context of the host-side hook.
Because the Linux dynamic linker processes \texttt{LD\_PRELOAD} variable before any application code runs \cite{linux_ldso_manpage}, the attacker's malicious shared library is loaded and executed with full root privileges on the host, achieving a complete container escape.
The significance of this vulnerability extends beyond its technical elegance. It strikes at a foundational assumption of container security: that isolation primitives reliably prevent a containerised process from affecting the host \cite{lin2018measurement}. In the context of multi-tenant AI cloud services ( where numerous organisations share the same physical GPU nodes ) a single exploitable container can compromise the entire host, potentially exposing sensitive model weights, training data, and API keys.

\newpage
\medskip
\noindent This paper is structured as follows:
\begin{itemize}
    \item Chapter \ref{sec:background} provides the necessary background on Linux container security, the NVIDIA Container Toolkit, the \texttt{LD\_PRELOAD} mechanism, and historically related vulnerabilities.
    \item Chapter \ref{sec:vuln} presents a deep technical analysis of CVE-2025-23266 itself.
    \item Chapter \ref{sec:exploitation} details the exploitation mechanics step by step.
    \item Chapter \ref{sec:impact} assesses the real-world impact.
    \item Chapter \ref{sec:mitigation} discusses mitigations and hardening strategies.
    \item Chapter \ref{sec:implications} reflects on broader security implications.
\end{itemize}    